\subsection*{Startup coefficients and the remaining source certification}
The fixed-radius inverse and two-space majorant reduce the analytic argument
to a common-domain problem, but they do not by themselves identify the actual
leading datum or verify the first two source equations. Here we audit those
equations against (5.2)--(5.6) of \cite{openai2026}.

\paragraph{Leading data and quantitative radius.}
Appendix B starts with a pressure-axis datum $\Pi_0(0,\eta)$ from its
Lemma A.5 and $U_*(\eta)=4\eta+j_0$, $0<j_0\le0.05$. Its choices
$\delta_*,\sigma_*>0$ and a simply connected neighborhood $\Omega$ obey
nonvanishing conditions for $L$ and $H_*^2+\sigma_*^2$. It sets
\[
\zeta_*=-\frac{LH_*}{H_*^2+\sigma_*^2},\qquad
\varphi_*(\eta)=\exp\left(\Lambda\int_0^\eta\zeta_*(w)\,dw\right),
\]
then chooses $\Lambda$ sufficiently large and $C$ sufficiently large
relative to the complex supremum of $\varphi_*$. Proposition B.2
provides a local analytic solution after those dependent choices.
No particular numerical values for all these objects and their
common complex suprema are supplied. Thus a numerical frozen-profile
certificate cannot be inferred from the existence construction alone.

\paragraph{Exact source startup.}
For $n\ge1$, put $\lambda_n=2hn$ and adopt the source's convention
that negative-index fields vanish. Write
\[
\Pi_n^{\rm cur}=2C^{-2}\int_0^X\varphi_0\varphi_n\,dx,\qquad
\Pi_n^{\rm low}=\int_0^X\left[
 C^{-2}\sum_{i=1}^{n-1}\varphi_i\varphi_{n-i}
 -\frac{\Omega_{n-1}}{2x}\right]dx.
\]
The positive-order gauge is $\Pi_n(0,\eta)=0$, whereas the
order-zero trace $\Pi_0(0,\eta)$ is prescribed separately. The
complete first two radial-momentum sources are
\[
\begin{aligned}
\Omega_0={}&T_{0,0}V_0+
 V_0(V'_0-V_0/(2X))+U_0Z_{0,0}V_0-2XV''_0,\\
\Omega_1={}&T_{0,1}V_1+
 V_0(V'_1-V_1/(2X))+V_1(V'_0-V_0/(2X))\\
 &+U_0Z_{0,1}V_1+U_1Z_{0,0}V_0
 -2XV''_1-Z^{[2]}_{0,0}V_0.
\end{aligned}
\]

\begin{proposition}[Exact first- and second-order forcing]
\label{prop:startup-source-check}
Let $\mathcal K_n$ be the four current-order blocks (47)--(50),
$\mathcal R_\varphi=\tfrac12J_2$, $\mathcal R_U=\tfrac12J_1$
and $w_n=(\varphi_n,U_n)^\top$. Then
\[
w_n=\mathcal K_nw_n+
\begin{pmatrix}\mathcal R_\varphi F_n^\varphi\\
\mathcal R_U F_n^U\end{pmatrix}\qquad(n=1,2),
\]
with exact lower-order sources
\begin{align*}
F_1^\varphi={}&-Z_{b,0}^{[2]}\varphi_0,\\
F_1^U={}&-Z_{c,0}^{[2]}U_0+
 Z_{-2A,1}\left[-\frac12\int_0^X\frac{\Omega_0}{x}\,dx\right],\\
F_2^\varphi={}&-Z_{b,1}^{[2]}\varphi_1+
 V_1(\varphi_1'+\varphi_1/X)+U_1Z_{b,1}\varphi_1,\\
F_2^U={}&-Z_{c,1}^{[2]}U_1+V_1U'_1+U_1Z_{c,1}U_1\\
&+Z_{-2A,2}\int_0^X
  \left(C^{-2}\varphi_1^2-\frac{\Omega_1}{2x}\right)dx .
\end{align*}
In particular $\Omega_0$ has no negative-index delayed viscosity,
$\Omega_1$ includes $-Z^{[2]}_{0,0}V_0$, and the current pressure
factor $2C^{-2}\varphi_0\varphi_n$ is not counted as lower forcing.
\end{proposition}
\begin{proof}
The ordered-pair set $i+j=1$ consists of $(0,1),(1,0)$ and
has no positive-positive pair. At $n=2$ the single positive-positive
pair is $(1,1)$. In both cases $\Omega_{n-1}$ is known at the
beginning of order $n$ and the source convention removes $V_{-1}$.
Separating these pairs in (25)--(28) gives the displayed terms.
The incompressibility formula (24) implies $V_j=O(X)$, so the
pressure integrands are axis-regular. Apply (34)--(35) with zero
positive-order axis integration constants.
\end{proof}

\paragraph{The physical-coordinate identity.}
With $q$ defined implicitly by
$\tau=1-t=q(1-\eta^2)$, $z=q^D\eta$, $X=r^2/(2q)$,
and $L=1-2h\eta^2$, $d=1-\eta^2$, exact differentiation at
fixed physical coordinates yields
\[
\begin{gathered}
q_t=-L^{-1},\quad q_z=2\eta q^{1-D}L^{-1},\quad
\eta_t=D\eta(qL)^{-1},\quad \eta_z=q^{-D}dL^{-1},\\
X_t=X(qL)^{-1},\qquad X_z=-2\eta Xq^{-D}L^{-1}.
\end{gathered}
\]
Thus $\partial_t(q^aF)=q^{a-1}T_aF$ and
$\partial_z(q^aF)=q^{a-D}Z_aF$ exactly, with $L>0$ for
$0<h<1/100$ and $|\eta|\le1$.

\begin{lemma}[Conditional physical residual reconstruction]
\label{lem:conditional-physical-residual}
Assume one fixed leading profile satisfies the order-zero inner
balances, every positive-order coefficient satisfies (23)--(28)
with the exact gauges, and the physical velocity/pressure series
and every derivative entering the Navier--Stokes residual converge
locally uniformly on compact preterminal subcores. Then the
reconstructed field is divergence-free and satisfies the unforced
Navier--Stokes equations pointwise there, including across the axis.
\end{lemma}
\begin{proof}
On each such compact set, $q>0$ and $L>0$ have uniform positive
lower bounds. The displayed exact chain rules permit termwise
physical differentiation. Absolute convergence and the assumed
derivative bounds justify product convolution. At every expansion
order, the physical divergence and momentum residual coefficients
are (23)--(28), and vanish by hypothesis; hence the residual
vanishes. Axis regularity follows from $V_n=O(X)$, the explicit
factor $r$ in $u_\theta$, and analyticity in $X=r^2/(2q)$.
\end{proof}

The source's own higher-order existence lemma permits
order-dependent analytic neighborhoods. Its subsequent smooth
summation uses order-dependent cutoffs and does not assert
convergence of the unmodified infinite series. In addition,
its whole-space theorem has smooth compactly supported forcing.
The outstanding unforced verification requires an explicit common
profile-radius certificate, genuine solutions of the startup
equations above, the complete higher-order convergence estimate,
and a source-specific physical residual check.
